% Part: first-order-logic
% Chapter: tableaux
% Section: soundness.tex

\documentclass[../../../include/open-logic-section]{subfiles}

\begin{document}

\iftag{FOL}
      {\olfileid{fol}{tab}{sou}}
      {\olfileid{pl}{tab}{sou}}
      
\olsection{Kasahihan}

\begin{explain}
Sistem !!^a{derivation}, kayata tableaux, iku \emph{sah} yen ora bisa
!!{derive} bab kang sejatine ora lumaku. Mula kasahihan iku sawijining
sipat aman kang dijamin kanggo sistem !!{derivation}. Gumantung marang
sipat teoretis-bukti kang dirembug, kita kepengin ngerti, upamane, yen
\begin{enumerate}
\item saben !!{derivable}~$!A$ iku \iftag{FOL}{valid}{tautologi};
\item yen !!a{sentence} bisa !!{derivable} saka ukara-ukara liyane,
  ukara iku uga dadi akibate;
\item yen sawijining himpunan !!{sentence}s ora konsisten, himpunan
  iku ora bisa dicukupi.
\end{enumerate}
Iki sipat-sipat wigati saka sistem !!a{derivation}. Yen salah sijine
ora lumaku, sistem !!{derivation} kasebut cacad---sistem iku bakal
!!{derive} kakehan. Mula, netepake kasahihan sistem !!a{derivation}
iku wigati banget.

Amarga kabeh sipat teoretis-bukti iki ditemtokake liwat !!{tableau}s
katutup kanthi sawijining wujud, mbuktekake (1)--(3) ing ndhuwur
mbutuhake bukti ngenani sipat semantis !!{tableau}s katutup. Kaping
pisan kita bakal netepake apa tegese !!a{signed
  formula} dicukupi ing sawijining struktur, banjur nuduhake yen, yen
!!{tableau} katutup,
ora ana struktur kang nyukupi kabeh asumsine. Sawise iku, (1)--(3)
dadi korolari saka asil iki.
\end{explain}

\begin{defn}
\iftag{FOL}{!!^a{structure}}{!!^a{valuation}}~$\iftag{FOL}{\Struct{M}}{\pAssign{v}}$
\emph{nyukupi} !!a{signed formula} $\sFmla{\True}{!A}$ yen lan mung
yen $\iftag{FOL}{\Sat{M}}{\pSat{v}}{!A}$, lan nyukupi
$\sFmla{\False}{!A}$ yen lan mung yen
$\iftag{FOL}{\Sat/{M}}{\pSat/{v}}{!A}$. $\iftag{FOL}{\Struct{M}}{\pAssign{v}}$
nyukupi himpunan !!{signed formula}s~$\Gamma$ yen lan mung yen nyukupi
saben $\sFmla{S}{!A} \in \Gamma$. $\Gamma$~iku \emph{bisa dicukupi}
yen ana \iftag{FOL}{!!a{structure}}{!!a{valuation}} kang nyukupi
himpunan iku, lan \emph{ora bisa dicukupi} yen ora ana.
\end{defn}

\begin{thm}[Kasahihan]
  \ollabel{thm:tableau-soundness}
  Yen $\Gamma$ nduweni !!{tableau} katutup, $\Gamma$ ora bisa dicukupi.
\end{thm}

\begin{proof}
Ayo kita ngarani sawijining cabang saka !!a{tableau} bisa dicukupi yen
lan mung yen himpunan !!{signed formula}s ing cabang iku bisa dicukupi,
lan ayo kita ngarani !!a{tableau} bisa dicukupi yen ngemot paling ora
siji cabang kang bisa dicukupi.

Kita nuduhake pratelan iki: ngambakake !!{tableau} kang bisa dicukupi
nganggo salah siji aturan inferensi mesthi ngasilake !!{tableau} kang
bisa dicukupi. Iki bakal mbuktekake teorema: saben !!{tableau} katutup
diasilake kanthi ngetrapake aturan inferensi marang !!{tableau} kang
mung dumadi saka asumsi-asumsi saka~$\Gamma$. Dadi, yen $\Gamma$ bisa
dicukupi, saben !!{tableau} kanggo himpunan iku uga bisa dicukupi.
Nanging !!{tableau} katutup cetha ora bisa dicukupi: saben cabang
ngemot $\sFmla{\True}{!A}$ uga $\sFmla{\False}{!A}$, lan ora ana
struktur kang bisa nyukupi lan ora nyukupi~$!A$ bebarengan.

Upamana kita duwe !!{tableau} kang bisa dicukupi, yaiku !!a{tableau}
kanthi paling ora siji cabang kang bisa dicukupi. Ngetrapake aturan
inferensi bisa nambahake !!{signed formula}s ing sawijining cabang
utawa mecah cabang dadi loro. Yen !!{tableau} nduweni cabang kang bisa
dicukupi lan cabang iku ora diambakake dening aplikasi aturan kang
dirembug, cabang kasebut tetep bisa dicukupi ing !!{tableau} kang wis
  diambakake; mula tableau kang wis diambakake bisa dicukupi. Dadi,
kita mung kudu nimbang kasus nalika aturan ditrapake marang cabang
kang bisa dicukupi.

Ayo $\Gamma$ dadi himpunan !!{signed formula}s ing cabang kasebut,
lan ayo $\sFmla{S}{!A} \in \Gamma$ dadi !!{signed formula} kang
dikenani aturan. Yen aturan kasebut ora mecah cabang, kita kudu
nuduhake yen cabang kang wis diambakake, yaiku $\Gamma$ bebarengan
karo kesimpulan aturan, isih bisa dicukupi. Yen aturan mecah cabang,
kita kudu nuduhake yen paling ora salah siji saka rong cabang asil
bisa dicukupi.
  
Kaping pisan, kita nimbang inferensi-inferensi kang ora mecah cabang.
\begin{enumerate}
\item Cabang diambakake kanthi ngetrapake $\TRule{\True}{\lnot}$
  marang $\sFmla{\True}{\lnot !B} \in \Gamma$. Banjur cabang kang wis
  diambakake ngemot !!{signed formula}s $\Gamma \cup
  \{\sFmla{\False}{!B}\}$. Upamana
  $\iftag{FOL}{\Sat{M}}{\pSat{v}}{\Gamma}$. Mligine,
  $\iftag{FOL}{\Sat{M}}{\pSat{v}}{\lnot !B}$. Mula,
  $\iftag{FOL}{\Sat/{M}}{\pSat/{v}}{!B}$, yaiku
  $\iftag{FOL}{\Struct{M}}{\pAssign{v}}$ nyukupi
  $\sFmla{\False}{!B}$.
\item Cabang diambakake kanthi ngetrapake $\TRule{\False}{\lnot}$
  marang $\sFmla{\False}{\lnot !B} \in \Gamma$: Gladhen.
\item Cabang diambakake kanthi ngetrapake $\TRule{\True}{\land}$
  marang $\sFmla{\True}{!B \land !C} \in \Gamma$, kang ngasilake
  rong !!{signed formula}s anyar ing cabang: $\sFmla{\True}{!B}$ lan
  $\sFmla{\True}{!C}$. Upamana
  $\iftag{FOL}{\Sat{M}}{\pSat{v}}{\Gamma}$, mligine
  $\iftag{FOL}{\Sat{M}}{\pSat{v}}{!B \land !C}$. Banjur
  $\iftag{FOL}{\Sat{M}}{\pSat{v}}{!B}$ lan
  $\iftag{FOL}{\Sat{M}}{\pSat{v}}{!C}$. Iki ateges
  $\iftag{FOL}{\Struct{M}}{\pAssign{v}}$ nyukupi
  $\sFmla{\True}{!B}$ lan $\sFmla{\True}{!C}$ kalorone.
\item Cabang diambakake kanthi ngetrapake $\TRule{\False}{\lor}$
  marang $\sFmla{\False}{!B \lor !C} \in \Gamma$: Gladhen.
\item Cabang diambakake kanthi ngetrapake $\TRule{\False}{\lif}$
  marang $\sFmla{\False}{!B \lif !C} \in \Gamma$: Iki ngasilake
  rong !!{signed formula}s anyar ing cabang: $\sFmla{\True}{!B}$ lan
  $\sFmla{\False}{!C}$. Upamana
  $\iftag{FOL}{\Sat{M}}{\pSat{v}}{\Gamma}$, mligine
  $\iftag{FOL}{\Sat/{M}}{\pSat/{v}}{!B \lif !C}$. Banjur
  $\iftag{FOL}{\Sat{M}}{\pSat{v}}{!B}$ lan
  $\iftag{FOL}{\Sat/{M}}{\pSat/{v}}{!C}$. Iki ateges
  $\iftag{FOL}{\Struct{M}}{\pAssign{v}}$ nyukupi
  $\sFmla{\True}{!B}$ lan $\sFmla{\False}{!C}$ kalorone.

\iftag{FOL}{%
\item Cabang diambakake kanthi ngetrapake $\TRule{\True}{\lforall}$
  marang $\sFmla{\True}{\lforall[x][!A(x)]} \in \Gamma$: Iki
  ngasilake !!{signed formula} anyar~$\sFmla{\True}{!A(t)}$ ing
  cabang. Upamana $\Sat{M}{\Gamma}$, mligine
  $\Sat{M}{\lforall[x][!A(x)]}$. Miturut
  \olref[syn][sem]{prop:quant-terms}, $\Sat{M}{!A(t)}$. Mula,
  $\Struct{M}$ nyukupi $\sFmla{\True}{!A(t)}$. Cathetan koreksi sumber
  OLPL-020: premis aturan Inggris nggunakake huruf formula kang beda
  saka asil lan kabeh argumentasi; wangun aturan umum netepake huruf
  formula kang dibalekake ing kene.

\item Cabang diambakake kanthi ngetrapake $\TRule{\False}{\lforall}$
  marang $\sFmla{\False}{\lforall[x][!A(x)]} \in \Gamma$: Iki
  ngasilake !!{signed formula} anyar~$\sFmla{\False}{!A(a)}$, ing
  ngendi $a$ iku !!a{constant} kang ora muncul ing~$\Gamma$. Amarga
  $\Gamma$ bisa dicukupi, ana $\Struct{M}$ saengga
  $\Sat{M}{\Gamma}$, mligine
  $\Sat/{M}{\lforall[x][!A(x)]}$. Kita kudu nuduhake yen $\Gamma \cup
  \{\sFmla{\False}{!A(a)}\}$ bisa dicukupi. Kanggo nindakake iki,
  kita nemtokake $\Struct{M'}$ kang trep kaya mangkene. Cathetan
  koreksi sumber OLPL-020 uga ditrapake ing kene: papat sebutan huruf
  formula liyane ing premis lan pratelan semantis dibalekake supaya
  cocog karo asil aturan lan kabeh langkah sabanjure.

  Miturut \olref[syn][ass]{prop:sat-quant},
  $\Sat/{M}{\lforall[x][!A(x)]}$ yen lan mung yen kanggo sawijining
  $s$, $\Sat/{M}{!A(x)}[s]$. Saiki ayo $\Struct{M'}$ padha karo
  $\Struct{M}$, kajaba $\Assign{a}{M'} = s(x)$. Miturut
  \olref[syn][ext]{cor:extensionality-sent}, kanggo saben
  $\sFmla{\True}{!C} \in \Gamma$, $\Sat{M'}{!C}$, lan kanggo saben
  $\sFmla{\False}{!C} \in \Gamma$, $\Sat/{M'}{!C}$, amarga $a$ ora
  muncul ing~$\Gamma$.

  Miturut \olref[syn][ext]{prop:extensionality},
  $\Sat/{M'}{!A(x)}[s]$. Miturut
  \olref[syn][ext]{prop:ext-formulas}, $\Sat/{M'}{!A(a)}[s]$. Amarga
  $!A(a)$ iku !!a{sentence}, miturut
  \olref[syn][ass]{prop:sentence-sat-true}, $\Sat/{M'}{!A(a)}$, yaiku
  $\Struct{M'}$ nyukupi $\sFmla{\False}{!A(a)}$.
\item Cabang diambakake kanthi ngetrapake $\TRule{\True}{\lexists}$
  marang $\sFmla{\True}{\lexists[x][!B(x)]} \in \Gamma$: Gladhen.
\item Cabang diambakake kanthi ngetrapake $\TRule{\False}{\lexists}$
  marang $\sFmla{\False}{\lexists[x][!B(x)]} \in \Gamma$: Gladhen.}{}
\end{enumerate}
Saiki ayo nimbang inferensi-inferensi kang mecah cabang.
\begin{enumerate}
\item Cabang diambakake kanthi ngetrapake $\TRule{\False}{\land}$
  marang $\sFmla{\False}{!B \land !C} \in \Gamma$, kang ngasilake
  rong cabang: cabang kiwa nerusake liwat $\sFmla{\False}{!B}$, lan
  cabang tengen liwat $\sFmla{\False}{!C}$. Upamana
  $\iftag{FOL}{\Sat{M}}{\pSat{v}}{\Gamma}$, mligine
  $\iftag{FOL}{\Sat/{M}}{\pSat/{v}}{!B \land !C}$. Banjur
  $\iftag{FOL}{\Sat/{M}}{\pSat/{v}}{!B}$ utawa
  $\iftag{FOL}{\Sat/{M}}{\pSat/{v}}{!C}$. Ing kasus kapisan,
  $\iftag{FOL}{\Struct{M}}{\pAssign{v}}$ nyukupi
  $\sFmla{\False}{!B}$, yaiku $\iftag{FOL}{\Struct{M}}{\pAssign{v}}$
  nyukupi formula-formula ing cabang kiwa. Ing kasus kapindho,
  $\iftag{FOL}{\Struct{M}}{\pAssign{v}}$ nyukupi
  $\sFmla{\False}{!C}$, yaiku $\iftag{FOL}{\Struct{M}}{\pAssign{v}}$
  nyukupi formula-formula ing cabang tengen.
\item Cabang diambakake kanthi ngetrapake $\TRule{\True}{\lor}$
  marang $\sFmla{\True}{!B \lor !C} \in \Gamma$: Gladhen.
\item Cabang diambakake kanthi ngetrapake $\TRule{\True}{\lif}$
  marang $\sFmla{\True}{!B \lif !C} \in \Gamma$: Gladhen.
\item Cabang diambakake nganggo \Cut: Iki ngasilake rong cabang,
  siji ngemot $\sFmla{\True}{!B}$ lan sijine ngemot
  $\sFmla{\False}{!B}$. Amarga
  $\iftag{FOL}{\Sat{M}}{\pSat{v}}{\Gamma}$ lan salah siji
  $\iftag{FOL}{\Sat{M}}{\pSat{v}}{!B}$ utawa
  $\iftag{FOL}{\Sat/{M}}{\pSat/{v}}{!B}$,
  $\iftag{FOL}{\Struct{M}}{\pAssign{v}}$ nyukupi cabang kiwa utawa
  cabang tengen.
\end{enumerate}
\end{proof}

\tagprob{FOL}
\begin{prob}
Jangkepana bukti \olref[fol][tab][sou]{thm:tableau-soundness}.
\end{prob}
\tagendprob

\tagprob{notFOL}
\begin{prob}
Jangkepana bukti \olref[pl][tab][sou]{thm:tableau-soundness}.
\end{prob}
\tagendprob

\begin{cor}
\ollabel{cor:weak-soundness}
Yen $\Proves !A$, mula $!A$ iku \iftag{FOL}{valid}{tautologi}.
\end{cor}

\begin{cor}
\ollabel{cor:entailment-soundness}
Yen $\Gamma \Proves !A$, mula $\Gamma \Entails !A$.
\end{cor}

\begin{proof}
  Yen $\Gamma \Proves !A$, mula kanggo sawetara $!B_1$, \dots,
  $!B_n \in \Gamma$, himpunan $\{\sFmla{\False}{!A},
  \sFmla{\True}{!B_1}, \dots, \sFmla{\True}{!B_n}\}$ nduweni
  !!{tableau} katutup. Miturut
  \olref{thm:tableau-soundness}, saben
  \iftag{FOL}{!!{structure}}{!!{valuation}}~$\iftag{FOL}{\Struct{M}}{\pAssign{v}}$
  nggawe sawetara $!B_i$ salah utawa nggawe $!A$ bener. Mula, yen
  $\iftag{FOL}{\Sat{M}}{\pSat{v}}{\Gamma}$, banjur uga
  $\iftag{FOL}{\Sat{M}}{\pSat{v}}{!A}$.
\end{proof}

\begin{cor}
\ollabel{cor:consistency-soundness}
Yen $\Gamma$ bisa dicukupi, mula himpunan iku konsisten.
\end{cor}

\begin{proof}
Kita mbuktekake kontraposisine. Upamana $\Gamma$ ora konsisten. Banjur
ana $!B_1$, \dots, $!B_n \in \Gamma$ lan !!{tableau} katutup kanggo
$\{\sFmla{\True}{!B_1}, \dots, \sFmla{\True}{!B_n}\}$. Miturut
\olref{thm:tableau-soundness}, ora ana
$\iftag{FOL}{\Struct{M}}{\pAssign{v}}$ saengga
$\iftag{FOL}{\Sat{M}}{\pSat{v}}{!B_i}$ kanggo kabeh $i=1$, \dots,~$n$.
Nanging banjur $\Gamma$ ora bisa dicukupi.
\end{proof}

\end{document}
